\section{Experiments}
\label{sec:experiments}

\subsection{Exploratory hardware evaluation}
\label{sec:setup}

The current measurements are development diagnostics. They do not constitute
the frozen final evaluation or establish search efficiency. TV denoising uses
$\lambda=0.1$ and three training plus three held-out $128\times128$ OCT
patches from the Kermany dataset's test split \citep{Kermany2018Identifying}.
Here held-out means excluded from recipe selection in that diagnostic.
Timing repetitions are not independent images, and no clinical denoising
efficacy claim follows from reaching an optimization tolerance.

The LP joint diagnostic uses two training and two held-out instances from the
previously frozen calibration partition, chosen for earlier solvability. The
instances are comp07-2idx and neos-1171737 for selection, with neos-1171448 and
neos-4300652-rahue excluded from selection. They are not a representative LP
sample. The original final LP holdout is untouched. Every measurement has three
warm repeats, a 20{,}000-iteration cap, and an independent original-coordinate
numerical check. The reports record preparation and cold-call costs separately.

\subsection{Execution-policy results}
\label{sec:results}

\begin{table}[t]
\centering\small
\begin{tabular}{llrlr}
\toprule
cell/device & tolerance & solves & selected policy & legacy/selected \\
\midrule
TV / CPU & $10^{-4}$ & 3/3 & float32, $K=50$ & $4.60\times$ \\
TV / CPU & $10^{-6}$ & 3/3 & float64, $K=200$ & $2.91\times$ \\
TV / RTX 4090 & $10^{-4}$ & 3/3 & schedule, $K=200$ & $5.20\times$ \\
TV / RTX 4090 & $10^{-6}$ & 3/3 & schedule, $K=200$ & $4.55\times$ \\
TV / RTX 5090 & $10^{-4}$ & 3/3 & schedule, $K=200$ & $4.37\times$ \\
TV / RTX 5090 & $10^{-6}$ & 3/3 & schedule, $K=200$ & $4.03\times$ \\
LP / RTX 4090 & $10^{-6}$ & 2/2 & XLA float64, $K=50$ & $2.25\times$ \\
LP / RTX 5090 & $10^{-6}$ & 2/2 & XLA float64, $K=50$ & $2.31\times$ \\
\bottomrule
\end{tabular}
\caption{Training-selected execution policies evaluated on small diagnostic
holdouts. Ratios compare the same mathematical solver with legacy execution
and the selected policy. They do not measure an advantage over competent
cadence/precision tuning or attribute a gain to LLM proposals.}
\label{tab:tv-results}
\end{table}

The TV algorithm is the previously banked Condat--V\~u policy
$(\tau,\sigma,\rho)=(0.05,1.0,1.9)$. Both GPUs select the same execution
policy; this gives no evidence of GPU-to-GPU specialization. The CPU-selected
recipe is slower on the GPUs in this diagnostic, but the reverse transfer is
not uniformly worse. Pure float32 loses some solves at $10^{-6}$; float64
refinement can restore the declared tolerance. Lower precision is consequently
an accuracy-dependent choice, not an unconditional speedup.

\subsection{Joint selection and kernel controls}

The LP catalog crosses two frozen parameter/relaxation configurations with
legacy execution, XLA cadence choices and three generated CUDA specializations.
Joint selection, algorithm-then-implementation, the reverse order and
implementation-only selection all choose frozen B0 with XLA float64 at $K=50$.
There is no demonstrated joint-selection advantage in this catalog. B0 uses
$(\tau,\sigma)=(1.98,0.495)$ with the frozen Ruiz/diagonal-preconditioning
configuration; it is not the full adaptive PDLP solver.

The thread-per-row CUDA variant at $K=50$ is about 25 percent slower than XLA
on the two held-out LP instances on the RTX 4090, and about 50 percent slower
on the RTX 5090. Warp variants also lose. On a separately traced trento1
diagnostic, XLA emits seven kernels per iteration while the generated thread
variant emits three, including loop bookkeeping. The serial dual-row update
dominates the generated variant. Fewer launches therefore do not establish
an end-to-end improvement. A cadence of 200 loses one of the two LP solves
under the iteration cap, illustrating the interaction between checking
frequency and detection of a transient threshold crossing.

\subsection{Search and verification evidence}

A corrected proposer pilot has two seeds and eight fresh evaluations per arm
on CPU, comparing independent random search, typed population mutation, and
LLM proposals with absent, identity-only or measured hardware context. It does
not establish an LLM or hardware-context benefit, and frequent fallbacks limit
interpretation. The earlier generation-4 versus generation-6 comparison used
coupled random streams and is excluded from sample-efficiency claims.

Historical verification auditing reports 285 successful numerical checks over
2{,}002 held-out evaluation rows; the remaining rows are logged gate rejections
or iteration-cap non-solves. These are repeated evaluations, not 2{,}002
independent instances. The OOD gate stress test rejects 1{,}340 configurations
under the supplied sufficient conditions; it does not establish algorithmic
non-transfer or divergence. A paired one-model repair experiment reports
52 percent success without feedback, 42 percent with generic feedback and
36 percent with exact-condition feedback within two retries. This is a scoped
negative observation, not a general limitation of LLM self-repair.

\subsection{Development factorial against a stronger baseline}

We cross two mathematically distinct TV updates, PDHG and Condat--V\~u, with
two parameter settings, three precision policies and three check cadences
($K=16,64,256$), yielding 36 recipes. A numerical map comparison confirms that
scheme labels correspond to different updates. The baseline is the banked
Condat--V\~u solver with compiled float64 execution and $K=64$. This improves
on the legacy denominator but does not replace full automatic configuration.
Twelve $256\times256$ OCT images from the dataset's training split provide six
selection and six development-validation images. A filename-derived audit
finds twelve distinct patient groups and no overlap between folds. The first
runner omitted those identifiers; the audit is recorded as a correction, and
the revised runner checks groups before timing. All three device runs use
identical input hashes and captured numerical source. We make no cross-device
throughput claim from these workstation diagnostics.

% Generated by scripts/summarize_solver_interaction.py from completed runs.
\begin{table}[t]
\centering\small
\begin{tabular}{llrrrr}
\toprule
device & tolerance & base (s) & impl. (s) & joint (s) & solves \\
\midrule
CPU & $10^{-4}$ & 0.09405 & 0.09405 & 0.09405 & 6/6/6 \\
CPU & $10^{-6}$ & 1.83258 & 1.82872 & 1.82872 & 5/5/5 \\
RTX 4090 & $10^{-4}$ & 0.01341 & 0.00911 & 0.00911 & 6/6/6 \\
RTX 4090 & $10^{-6}$ & 1.31134 & 1.28160 & 1.28160 & 5/5/5 \\
RTX 5090 & $10^{-4}$ & 0.01466 & 0.00933 & 0.00933 & 6/6/6 \\
RTX 5090 & $10^{-6}$ & 1.31495 & 1.30005 & 1.29043 & 5/5/5 \\
\bottomrule
\end{tabular}
\caption{Larger-patch development diagnostic: six training and six validation OCT images at $256\times256$. Entries are penalized SGM-1 seconds, with a 120-second non-solve penalty and 5,000-iteration cap. Solves lists baseline/implementation-only/joint counts, each out of six. The baseline is compiled banked Condat--V\~u with float64 and $K=64$. Selections use training data only. This is a finite catalog, not a matched-budget search experiment.}
\label{tab:development-interaction}
\end{table}


At $10^{-4}$, implementation selection reduces validation SGM-1 from 13.41
to 9.11\,ms on the RTX 4090 ($1.47\times$) and from 14.66 to 9.33\,ms on
the RTX 5090 ($1.57\times$), with all six images solved. Both choose the
banked Condat--V\~u solver with float32-to-float64 refinement and $K=256$;
CPU selection retains the baseline. Joint selection matches
implementation-then-algorithm selection on both GPUs. Algorithm-first
selection is about three percent slower on the RTX 4090 and ties on the
RTX 5090. The data therefore do not establish a benefit over both sequential
orders.

At $10^{-6}$, all selected recipes solve five of six validation images under
the 5,000-iteration cap. The fixed 120-second failure penalty dominates
the reported scores; these are not mean successful-solve times. The RTX 5090
joint choice differs from the sequential choices but improves their validation
penalized score by only about 0.75 percent, following a very small training
score difference. This single diagnostic does not establish a robust joint
advantage. Pure float32 Condat--V\~u solves none of the six at this tolerance.
Native TV kernels are absent from this catalog.

The main submission comparison still requires matched-budget search arms,
larger independent validation sets, a measured hardware cost explanation and
a frozen final evaluation. These catalog selections are not independent
evolutionary runs. A second mathematical problem family remains necessary
for the intended breadth of claims.
